% Compile with pdfLaTeX.
\documentclass[10pt,a4paper]{article}

\usepackage[T1]{fontenc}
\usepackage[utf8]{inputenc}
% Palatino and Helvetica are the pdfLaTeX counterparts closest to the
% TeX Gyre Pagella / Adventor pairing used in the original CV.
\usepackage{mathpazo}
\usepackage[scaled=0.92]{helvet}

\usepackage[a4paper,left=17mm,right=17mm,top=14mm,bottom=14mm]{geometry}
\usepackage{xcolor}
\usepackage{enumitem}
\usepackage{tabularx}
\usepackage{titlesec}
\usepackage{hyperref}
\usepackage{microtype}
\usepackage{needspace}

\definecolor{accent}{RGB}{35,132,158}
\definecolor{textgray}{RGB}{82,86,89}
\definecolor{lightgray}{RGB}{150,153,155}

\hypersetup{
  colorlinks=true,
  urlcolor=accent,
  linkcolor=accent,
  pdfauthor={Nikolay Tulupov},
  pdftitle={Nikolay Tulupov -- Radio Systems and Electronics Engineer}
}

\pagestyle{empty}
\setlength{\parindent}{0pt}
\setlength{\parskip}{0pt}
\setlength{\tabcolsep}{0pt}
\linespread{1.08}
\emergencystretch=2em

\titleformat{\section}
  {\needspace{4\baselineskip}\sffamily\bfseries\large\color{accent}}
  {}{0pt}{\MakeUppercase}
  [\vspace{0.10em}\color{accent}\titlerule]
\titlespacing*{\section}{0pt}{1.00em}{0.55em}

\setlist[itemize]{
  leftmargin=1.25em,
  itemsep=0.28em,
  topsep=0.32em,
  parsep=0pt,
  partopsep=0pt
}

\newcommand{\cvrole}[4]{%
  \begin{tabularx}{\textwidth}{@{}>{\raggedright\arraybackslash}X
    >{\raggedleft\arraybackslash}p{0.27\textwidth}@{}}
    {\sffamily\bfseries #1} & {\small\color{textgray}#4} \\
    {\itshape #2} & {\small\color{lightgray}#3}
  \end{tabularx}
  \vspace{0.12em}
}

\newcommand{\project}[2]{%
  \needspace{3\baselineskip}
  {\sffamily\bfseries #1}\par
  \vspace{0.12em}
  #2
}

\newcommand{\skillrow}[2]{%
  \textbf{#1}\hspace{0.45em}#2\par
  \vspace{0.13em}
}

\begin{document}

% Header
\begin{center}
  {\fontsize{24}{27}\selectfont\sffamily\bfseries\color{accent}
   NIKOLAY TULUPOV}\par
  \vspace{0.22em}
  {\large\sffamily\color{textgray}
   Radio Systems \& Electronics Engineer}\par
  \vspace{0.48em}
  {\small
    Moscow, Russia
    \enspace\textcolor{lightgray}{|}\enspace
    \href{tel:+79307297552}{+7 930 729-75-52}
    \enspace\textcolor{lightgray}{|}\enspace
    \href{mailto:tulupov.nd@phystech.edu}{tulupov.nd@phystech.edu}
    \enspace\textcolor{lightgray}{|}\enspace
    \href{https://t.me/oncoming_lane}{Telegram}
  }\par
  \vspace{0.16em}
  {\small
    \href{https://github.com/oncominglane}{GitHub}
    \enspace\textcolor{lightgray}{|}\enspace
    \href{https://www.researchgate.net/profile/Nikolay-Tulupov}{ResearchGate}
    \enspace\textcolor{lightgray}{|}\enspace
    \href{https://orcid.org/0009-0004-2017-2288}{ORCID: 0009-0004-2017-2288}
  }
\end{center}

\section{Profile}

Radio systems and electronics engineer with over two years of experience in
end-to-end development of hardware--software systems. Designed a remote-control
system for a Motorola XTL2500 radio, covering circuit and PCB design, embedded
interfaces, the SB9600 protocol, real-time bidirectional audio, Linux software,
system integration, testing, and technical documentation. Experienced with
radio equipment, measurement instruments, single-board computers, and
microcontrollers.

\section{Core Expertise}

\skillrow{Radio \& Electronics:}{radio equipment, circuit design, schematic
capture, PCB design, prototyping, hardware debugging, test instruments}
\skillrow{Embedded \& Interfaces:}{microcontrollers, single-board computers,
UART, SPI, I2C, RS-232, GPIO, ADC/DAC, CAN}
\skillrow{Software \& Systems:}{C/C++, Python, Linux, CMake, TCP/IP,
client--server systems, real-time audio, Git}
\skillrow{Engineering Tools:}{EasyEDA, MATLAB/Simulink, Motor-CAD, Xilinx ISE,
oscilloscope, soldering equipment}

\section{Relevant Experience}

\cvrole{TIENDA LLC}
       {Engineer, Technical Department}
       {Moscow, Russia}{Apr 2026 -- Present}
\begin{itemize}
  \item Continue development of the Motorola XTL2500 remote-control system
        following the project's organizational transfer from PromSvyazRadio.
  \item Develop and debug its hardware and software components and maintain
        project technical documentation.
\end{itemize}

\vspace{0.35em}
\cvrole{PromSvyazRadio LLC}
       {Assistant to the Chief Engineer, Technical Department}
       {Moscow, Russia}{Feb 2024 -- Apr 2026}
\begin{itemize}
  \item Developed a distributed client--server system on Linux-based
        single-board computers for remote control of a Motorola XTL2500 radio.
  \item Implemented real-time bidirectional audio streaming, command handling,
        and interaction with the radio through the SB9600 protocol.
  \item Parsed SB9600 frames and radio-display data and integrated UART, SPI,
        I2C, ADC/DAC, and network interfaces.
  \item Designed the system schematic and PCB and contributed to the enclosure,
        assembly, testing, and technical documentation.
  \item Performed end-to-end hardware/software integration and debugged the
        system using laboratory measurement equipment.
\end{itemize}

\vspace{0.35em}
\cvrole{AmpereMagnete}
       {Engineering Intern, Artificial Intelligence Group}
       {Troitsk, Moscow, Russia}{Mar 2025 -- Mar 2026}
\begin{itemize}
  \item Developed embedded firmware for an AT32-based power inverter and
        software for remote inverter control and test-bench operation.
  \item Implemented CAN-based telemetry acquisition, visualization, logging,
        and experimental data analysis.
  \item Co-developed StandControl, a registered client--server application for
        remote control and monitoring of a power inverter.
\end{itemize}

\newpage

\section{Additional Engineering Experience}

\cvrole{Moscow Institute of Physics and Technology}
       {Research Engineer, Robotics and Electric Drive Laboratory}
       {Moscow, Russia}{Mar 2026 -- Present}
\begin{itemize}
  \item Develop software for experimental test-bench operation, data
        processing, visualization, and automated engineering calculations.
  \item Model electric-drive systems and experimentally validate control
        algorithms on physical equipment.
  \item Automate Motor-CAD calculations and electric-motor geometry optimization.
\end{itemize}

\section{Selected Engineering Projects}

\project{Remote-Control System for the Motorola XTL2500 Radio}{
  \begin{itemize}
    \item Designed an end-to-end hardware--software system for operating a radio
          remotely over an IP network.
    \item Developed C/C++ software for Linux-based single-board computers,
          including client--server communication and real-time bidirectional
          audio using ALSA.
    \item Implemented SB9600 command exchange, frame parsing, display decoding,
          and integration with UART, SPI, I2C, ADC/DAC, and Ethernet interfaces.
    \item Designed the schematic and PCB and participated in enclosure design,
          system assembly, hardware debugging, and documentation.
  \end{itemize}
}

\vspace{0.25em}
\project{StandControl: Power-Inverter Control and Monitoring System}{
  \begin{itemize}
    \item Co-developed a client--server application for remote inverter control
          over CAN, including torque, speed, and current operating modes.
    \item Implemented telemetry monitoring, visualization, logging, network
          communication, and error handling using Python and C/C++.
    \item Registered as a computer program in 2026
          (\href{https://new.fips.ru/registers-doc-view/fips_servlet?DB=EVM\&DocNumber=2026614166\&TypeFile=html}{No.~2026614166}).
  \end{itemize}
}

\vspace{0.25em}
\project{OptiRunMotorCAD}{
  Developed a software application for automating Motor-CAD calculations and
  electric-motor geometry optimization. Registered as a computer program in
  2025
  (\href{https://new.fips.ru/registers-doc-view/fips_servlet?DB=EVM\&DocNumber=2025693472\&TypeFile=html}{No.~2025693472}).
}

\section{Education}

\cvrole{Moscow Institute of Physics and Technology}
       {M.Sc. in Applied Mathematics and Physics}
       {Moscow, Russia}{2026 -- 2028 (expected)}
{\small\textit{Specialization: E-Mobility -- Technologies, Materials and Control}}\par
\vspace{0.25em}
\cvrole{Moscow Institute of Physics and Technology}
       {B.Sc. in Applied Mathematics and Physics}
       {Moscow, Russia}{2022 -- 2026}
{\small\textit{Concentration: Radio Engineering and Computer Technology}}\par

\section{Selected Publications \& Recognition}

\begin{itemize}
  \item N. Tulupov and D. Kanurin,
        \href{https://doi.org/10.1109/EDM69524.2026.11632234}{``Physics-Informed TD3 for End-to-End Torque Control of a PMSM Drive,''}
        \textit{2026 IEEE 27th International Conference of Young Professionals
        in Electron Devices and Materials (EDM)}, 2026.
  \item Best Presentation Award, Exact and Engineering Sciences section,
        III International Youth Conference ``AI Technologies in Science and
        Education,'' Moscow State University, 2025.
\end{itemize}

\section{Languages}

Russian -- Native \hspace{1.4em}\textcolor{lightgray}{|}\hspace{1.4em}
English -- B2 (Upper-Intermediate)

\end{document}
